\documentclass[12pt]{article}

\usepackage[top=1in, bottom=1in, left=1in, right=1in]{geometry}
\usepackage{parskip}
\usepackage[colorlinks=true,linkcolor=blue,citecolor=Mahogany,urlcolor=blue]{hyperref}

\usepackage[usenames, dvipsnames]{xcolor}
\newcommand{\blue}[1]{\textcolor{blue}{#1}}

\begin{document}

\section*{Possible questions}

\tableofcontents

\section{Core Qs - preferences for redistribution}

    \blue{Shifting the tax cutoff as a treatment}
    
    Imagine that the government is considering raising the tax rate on incomes above \$500,000 by 10\% (from 35-37\% to 45-47\%) and decreasing the tax rate on incomes under \$75,000 by 10\% (from 12-22\% to 2-12\%). Would you support this policy?
    % Assume that the government's budget would be unchanged.

    \blue{Testing for ingroup favoritism}
    
    Using the strategy from \cite{fischbacher2023identity}, give participants information about a hypothetical person of a given race, political orientation, nationality, income level, etc. Tell them that the person's income was earned, random, or unfair. Ask them to take income from or give income to this person from another hypothetical person.

    \blue{open-ended questions?}

\section{Core Qs - economic effects of redistribution}

    \blue{In response to a redistribution question:}
    Would this policy reduce or increase the growth of the U.S. economy?

    Would this policy reduce or increase entrepreneurship?

\section{Core Qs - beliefs about mobility}

    Most people in the United States who want to get ahead can make it if they are willing to work hard.

    Under a fair economic system, people with higher abilities would earn higher salaries.

    Competition is good. It stimulates people to work hard and develop new ideas.

    Wealth can grow so there’s enough for everyone.

    Are poorer people generally lazier?

    Do poor people deserve less money?

\section{Miscellaneous Qs}

    What is the main reason you support/oppose this tax policy?
    \begin{itemize}
        \item It would save/lose me money
        \item It would be good/bad for the economy
        \item It is fair/unfair
        \item It helps/hurts people in need
    \end{itemize}

    Did you vote/who for.

\section{Consumption/income preferences}

    \subsection{Uncertain future income/min max}

        \blue{Ask how confident they are about their future income path. Those less certain should prefer redistribution more.}

        How much do you expect to earn this year in U.S. dollars, before taxes?

        What do you expect your annual income to be 5 years from now?

        What do you expect your annual income to be at the peak of your career?

        What is the lowest you might make?

        What is the highest you might make?

        How certain of this are you? How likely is this?

        What is the likelihood that you will earn within \$10,000 of this amount?

        \blue{Gauge their risk tolerance. Given a level of income uncertainty, those with less risk tolerance should want more redistribution. Risk tolerance is a good covariate to include because those with higher risk aversion might want more redistribution all else equal, but they'll also choose safer career paths.}
    
    \subsection{Maximizing current income/high discount rate}

    \blue{Ask about their expected future income. Young people who earn the same amount in the near-term should have similar policy responses even if their income paths diverge in the long-term.}

    \blue{Gauge their intertemporal preferences directly by asking if they'd prefer x amount now or x amount later. Or if they'd want a temporary income cut now or later. See if this correlates with policy responses.}
    
    \subsection{More/less materialistic}

    \blue{If young people are less materialistic, they may be more willing to give up money for another cause/accept a tax if the experimental cutoff hits them.}

    How important is making money to you?

    When you were younger, how important was making money to you?

    How much more content would you be if you made more money?

    Economic growth will make us happier.

    \subsection{More/less career focused}
    
    Is it impressive to have a lot of money?
    
    Do people in the U.S. find those who have a lot of money impressive?

    Do you want to get out of your current social class/income level?

\section{Leisure preferences}

    \subsection{The young find hard work painful}

    \blue{Are there any long-running questions in the WVS/ANES/GSS?}

    How many hours a week are you willing to work to [insert desirable outcome]?

    How much free time would you sacrifice to earn [x amount of money] more this year?

    How important is work-life balance to you? 

    How many hours a week do you work? / How many hours a week did you work in your 20s? 

    \subsection{The young value activities outside of work}

\section{Value preferences}
    
    \subsection{Individualism}

        \blue{Haidt's Moral Values Questionnaire has questions}

        \blue{You could ask a ``control'' redistribution question that doesn't involve politics/economic reasoning but does gauge their generosity and utility coming from others' success.}
    
    \subsection{Minority vs. majority/ingroup vs. outgroup}

        \blue{Find questions from literature}
    
    \subsection{Universalist vs. communalist}

        \blue{Haidt's Moral Values Questionnaire has questions}
    
    \subsection{Openness to experience}

        \blue{Find questions from literature}

\section{Beliefs about incentive effects}
    
    \subsection{Family's mobility}

    What social class were your parents when they were children?

    What is your parent or guardians’ current combined annual income?

    Did your parents “move up” in their income level or social class over their lifetime? Did they improve their economic position? How much?
    
    \subsection{Media/societal portrayal of mobility}

    \blue{Ask them what the media sentiment is about income mobility right now, or what the average person believes. Or, what people on the left and right believe separately.}

    Do people in the United States today believe it's possible to move up with hard work?

    When you were young, did people in the United States believe it's possible to move up with hard work?
    
    \subsection{Temporary optimism/belief in government}

    \blue{Older people may believe less in government and have a narrower view of ``what will work.'' Older people should be able to recall more instances where the government has failed them, and the amount of instances may be correlated with their policy preferences.}

    It’s often beneficial when the government steps in to solve problems.
    \textit{True; False}

    \blue{You could ask a question where the government's budget is fixed and the proposed policy purely redistributed the tax burden, versus a question where there is a proposed tax on the rich that funds additional government programs.}

\section{Behavioral biases}
    
    \subsection{The young don't know how much income they'll need/want in the future}

    \blue{Ask young people to estimate how much money they'll need to live comfortably in the future. Ask old people how much they'll need to live comfortably right now. See if the latter answer is higher.}

    How much income would you need to live comfortably now?

    How much income will you need to live comfortably when you are [insert age]?
    
    \subsection{The young don't anticipate how much taxes will hurt}

    \blue{Ask a question related to redistribution or monetary reward/punishment that doesn't involve taxes. See if young/old people who answer similarly on this question answer differently on the tax redistribution question.}

    \blue{One argument is that young people these days age slower/have adult responsibilities at a later point. If their experiences with taxation/government matter, then this could be making today's young more liberal (than older generations when they were the same age).}

\section{Demographics}

    What is your age?

    How many children do you have?

\newpage
\bibliographystyle{apalike}
\bibliography{refs}

\end{document}